\chapter{Planteamiento del problema}

\section{Antecedentes}
Esta sección presenta una revisión de la literatura académica relacionada con la arquitectura de software, su relación con el consumo de recursos computacionales y la eficiencia energética del software. La búsqueda se orientó siguiendo los principios generales de la metodología PRISMA (Preferred Reporting Items for Systematic Reviews and Meta-Analyses), cuyo proceso se documenta de manera transparente en el apartado 5.7. 

\subsection{Arquitectura de software y consumo de recursos}
La relación entre las decisiones de diseño arquitectónico y el consumo de recursos computacionales ha sido objeto de estudio desde hace más de una década. Seo et al. [11] propusieron un marco para estimar el impacto del estilo arquitectónico de un sistema distribuido sobre su consumo energético, sentando una de las primeras bases metodológicas para abordar este tipo de análisis de forma sistemática. Años más tarde, Jagroep et al. [6] extendieron las vistas de arquitectura de software tradicionales para incorporar una perspectiva de consumo energético, mediante un caso de estudio sobre el consumo de recursos en software empresarial, evidenciando que las decisiones arquitectónicas tienen efectos medibles sobre el uso de CPU y memoria. Estos trabajos constituyen unantecedente metodológico directo para la presente investigación, en la medida en que demuestran la viabilidad de evaluar experimentalmente el consumo de recursos asociado a distintas decisiones de diseño arquitectónico.

\subsection{Arquitecturas monolíticas y modularidad} 
Lacomparaciónentrearquitecturas monolíticas y modulares–particularmente bajo el paradigma de microservicios– ha sido ampliamente documentada desde la perspectiva de la escalabilidad y la mantenibilidad, y más recientemente desde la perspectiva del consumo de recursos. AlDebagy y Martinek [3] realizaron una revisión comparativa de las arquitecturas de microservicios y monolíticas, identificando sus principales diferencias estructurales y operativas. Por su parte, Su, Li y Taibi [4] desarrollaron una revisión multivocal sobre el fenómeno inverso–la migración de microservicios hacia arquitecturas monolíticas–, evidenciando que la decisión arquitectónica no es unidireccional y que existen contextos en los que la modularidad introduce una complejidad que no se traduce en beneficios proporcionales. Este antecedente resulta relevante porque cuestiona la idea de que la modularidad es siempre preferible, y refuerza la necesidad de evaluar caso por caso el impacto de cada estilo arquitectónico, incluyendo su impacto energético. 5.3 Modularidad, rendimiento y consumo energético Unconjuntomásespecíficodeestudioshaabordadodirectamentelacomparacióndelconsumo energético entre arquitecturas monolíticas y modulares. Capuano, O’Dea y Muccini [5] realizaron unanálisis comparativo experimental entre versiones monolíticas y de microservicios de aplicaciones de referencia (PetClinic y TicketMonster), encontrando que las arquitecturas de microservicios pueden consumir en promedio un 5.02\% menos de energía que sus contrapartes monolíticas bajo cargas medias yaltas, atribuyendo esta diferencia a una mejor utilización de recursos y a la ejecución modular de los componentes. Sin embargo, los propios autores advierten que estos resultados se limitan a condiciones controladas y a un número reducido de casos de estudio, por lo que no deben generalizarse sin evidencia adicional. En contraste, Araújo et al. [2], en una revisión sistemática de la literatura sobre consumo energético en arquitecturas de microservicios, señalan que los hallazgos disponibles no son consistentes entre sí y dependen fuertemente de las condiciones experimentales–tipo de aplicación, infraestructura de despliegue y métricas utilizadas–, lo que dificulta establecer conclusiones generalizables sobre si la modularidad reduce o incrementa el consumo energético. Esta revisión constituye uno de los antecedentes más directamente relacionados con el problema de investigación planteado, pues evidencia de manera explícita el vacío de conocimiento que esta investigación busca abordar.

\subsection{Consumoenergético y eficiencia energética del software}
Desde una perspectiva más amplia, el campo de la ingeniería de software sostenible ha generado un cuerpo de conocimiento relevante para contextualizar el problema de investigación. Capra, Francalanci y Slaughter [7] examinaron la relación entre los entornos de desarrollo de software y la eficiencia energética en aplicaciones de código abierto, aportando una de las primeras aproximaciones empíricas a la llamada "sostenibilidad del software". Kern et al. [1] propusieron criterios de evaluación para la eficiencia de recursos y energía en productos de software sostenible, ofreciendo un marco conceptual útil para definir métricas de consumo energético aplicables a este estudio. En una línea similar, Kipp, Jiang y Fugini [12] introdujeron un conjunto de métricas verdes para sistemas de tecnologías de la información conscientes de la energía, las cuales fueron posteriormente ampliadas por Kipp et al. [13] mediante la propuesta de indicadores de desempeño verde organizados por capas, aplicables a distintos niveles de la arquitectura de un sistema. Khanet al. [8] realizaron un mapeo sistemático de técnicas de software orientadas a mejorar la eficiencia energética de los centros de datos en la nube, identificando que buena parte de las estrategias disponibles se concentran en la infraestructura y no en decisiones de diseño arquitectónico a nivel de aplicación, lo cual refuerza la pertinencia de estudiar específicamente el efecto del estilo arquitectónico del software sobre el consumo energético. 

\subsection{Estudios de mapeo y revisión sobre software sostenible}
Anwar y Pfahl [10] propusieron el uso de analítica de software para avanzar hacia una ingeniería de software más sostenible, mientras que Marimuthu y Chandrasekaran [9] realizaron un mapeo sistemático de los aspectos de la ingeniería de software relacionados con la sostenibilidad y el software verde, identificando una notable escasez de estudios experimentales que relacionen directamente decisiones arquitectónicas específicas–como la elección entre monolito y microserviciosconmétricas de consumoenergéticomedidasempíricamente. Estehallazgo constituye un antecedente clave que justifica la necesidad de un estudio experimental comparativo como el que se propone en este proyecto.

\subsection{Vacío identificado}
La revisión de los antecedentes anteriores permite identificar un vacío de conocimiento consistente: si bien existe evidencia de que el estilo arquitectónico de un sistema de software influye en su consumo de recursos y energía [6,11], y existen estudios experimentales puntuales que comparan arquitecturas monolíticas y modulares en este aspecto [5], los resultados disponibles no son consistentes entre sí [2] y dependen en gran medida de las condiciones experimentales específicas de cada estudio. Adicionalmente, la mayoría de los marcos conceptuales sobre software sostenible [1,12,13] se han desarrollado de forma independiente de los estudios comparativos entre estilos arquitectónicos, sin que exista una integración clara entre ambos cuerpos de conocimiento. Este vacío–la ausencia de evidencia experimental consistente y reproducible sobre las condiciones bajo las cuales una arquitectura modular resulta más o menos eficiente energéticamente que una arquitectura monolítica equivalente– constituye el fundamento sobre el cual se plantea el problema de investigación de este proyecto. 

\subsection{Metodología PRISMA aplicada a la búsqueda de antecedentes}
La identificación de los antecedentes presentados en este capítulo se orientó siguiendo los principios de transparencia y documentación de la metodología PRISMA. A continuación se documenta el proceso realizado, indicando explícitamente los elementos que corresponden a búsquedas efectivamente realizadas y aquellos que quedan pendientes de una revisión sistemática formal en una etapa posterior del proyecto. Bases de datos consultadas: motor de búsqueda académico general (equivalente a Google Scholar), con verificación cruzada en repositorios de editoriales académicas (IEEE Xplore, ACM Digital Library, ScienceDirect, Springer Link, arXiv). Términosdebúsquedaempleados: "energyconsumptionmicroservicesmonolithicarchitecture"; "green software engineering energy efficiency systematic review"; "modularity software architecture performance trade-off". Periodo debúsqueda: sinrestricción de fecha inicial; se priorizaron publicaciones recientes sin excluir trabajos clásicos que definen conceptos fundamentales del área (por ejemplo, [11], de 2008). Criterios de inclusión: publicaciones académicas revisadas por pares (artículos de revista, ponencias de conferencia, capítulos de libro académico); relación directa con arquitectura de software, modularidad, consumo de recursos o eficiencia energética del software. Criterios de exclusión: entradas de blog, documentación comercial no académica, contenido sin autoría verificable, fuentes que abordan fenómenos ajenos al objeto de estudio (por ejemplo, eficiencia energética de hardware sin relación con decisiones de software). Identificación, cribado, elegibilidad e inclusión: el proceso de identificación y cribado se realizó de manera exploratoria mediante las búsquedas señaladas; el número exacto de resultados totales arrojados por cada motor de búsqueda, el número de duplicados eliminados y el número de documentos examinados en su totalidad no fue registrado de manera sistemática en esta fase preliminar del proyecto. [DATOPRISMAPORCOMPLETAR:cifrasexactasdeidentificación,cribado,elegibilidad e inclusión, a documentar en la fase de revisión sistemática formal del estado del arte]. De los documentos identificados, se incluyeron finalmente 13 referencias que cumplieron los criterios de inclusión y que sustentan directamente los distintos componentes del problema de investigación planteado, según se detalla en la tabla 5.1.

\subsection{Tabla de antecedentes}

\begin{table}[htbp]
\centering
\caption{Síntesis de antecedentes académicos revisados}
\label{tab:antecedentes}
\small
\begin{tabular}{|c|l|p{3cm}|p{3cm}|p{2.5cm}|p{3.5cm}|p{3cm}|}
\hline
\textbf{N.°} & \textbf{Autor(es) y año} & \textbf{Título} & \textbf{Objetivo} & \textbf{Metodología} & \textbf{Principales resultados} & \textbf{Relación} \\ \hline
1 & Seo et al. (2008) [11] & Marco para estimar el impacto arquitectónico en el consumo energético & Proponer un método de estimación del consumo energético según el estilo arquitectónico & Propuesta de marco conceptual y experimental & Sienta las bases metodológicas para relacionar arquitectura y energía & Antecedente metodológico directo \\ \hline
2 & Jagroep et al. (2017) [6] & Extensión de vistas de arquitectura con perspectiva energética & Incorporar el consumo energético como vista arquitectónica & Caso de estudio en software empresarial & Las decisiones arquitectónicas afectan mediblemente el consumo de recursos & Sustenta la relación entre arquitectura y recursos \\ \hline
3 & Al-Debagy y Martinek (2018) [3] & Revisión comparativa de microservicios y arquitecturas monolíticas & Comparar estructuralmente ambos estilos & Revisión comparativa & Identifica diferencias estructurales y operativas & Fundamenta la distinción de estilos (monolito vs. modular) \\ \hline
4 & Su, Li y Taibi (2024) [4] & De microservicio a monolito: revisión multivocal & Analizar la migración inversa (microservicio a monolito) & Revisión multivocal de literatura & La modularidad no siempre reporta beneficios proporcionales & Cuestiona la superioridad universal de la modularidad \\ \hline
5 & Capuano, O'Dea y Muccini (2026) [5] & Análisis comparativo de consumo energético: monolito VS. microservicios & Comparar el consumo energético entre ambos estilos & Experimento controlado (PetClinic, TicketMonster) & Ahorro energético promedio de 5.02\% con microservicios bajo carga media/alta & Evidencia experimental directa del problema \\ \hline
6 & Araújo et al. (2024) [2] & Consumo energético en arquitecturas de microservicios: revisión sistemática & Sintetizar la evidencia disponible sobre consumo energético en microservicios & Revisión sistemática de literatura (IEEE Access) & Resultados inconsistentes según condiciones experimentales & Evidencia directa del vacío de conocimiento \\ \hline
7 & Capra, Francalanci y Slaughter (2012) [7] & ¿Es "verde" el software? & Evaluar la eficiencia energética de entornos de desarrollo & Estudio empírico en software libre & Aporta una de las primeras aproximaciones a la sostenibilidad del software & Fundamenta el área de estudio (sostenibilidad) \\ \hline
8 & Kern et al. (2018) [1] & Productos de software sostenibles: criterios de evaluación & Proponer criterios de evaluación de eficiencia energética & Propuesta conceptual / marco de criterios & Define criterios de recursos y energía para productos de software & Aporta el marco conceptual de eficiencia energética \\ \hline
9 & Kipp, Jiang y Fugini (2011) [12] & Métricas verdes para sistemas de TI conscientes de la energía & Proponer métricas de consumo energético en sistemas de TI & Propuesta de métricas & Introduce indicadores de consumo energético aplicables a software & Sustenta la medición de métricas energéticas \\ \hline
10 & Kipp et al. (2012) [13] & Indicadores de desempeño verde por capas & Ampliar las métricas verdes a una estructura por capas & Propuesta de indicadores estratificados & Permite evaluar el consumo energético en distintos niveles arquitectónicos & Sustenta la evaluación energética por componentes \\ \hline
11 & Khan et al. (2020) [8] & Técnicas de software para centros de datos energéticamente eficientes & Mapear técnicas de eficiencia energética en la nube & Mapeo sistemático & Las estrategias existentes se centran en infraestructura, no en diseño de aplicación & Evidencia un vacío sobre la relación arquitectura-software \\ \hline
12 & Anwar y Pfahl (2017) [10] & Hacia una ingeniería de software más verde mediante analítica de software & Explorar el uso de analítica para la sostenibilidad del software & Mapeo sistemático & Identifica oportunidades de aplicar analítica al software sostenible & Aporta antecedentes metodológicos / complementarios \\ \hline
13 & Marimuthu y Chandrasekaran (2017) [9] & Aspectos de ingeniería de software verde y sostenible & Mapear el estado del arte en software verde y sostenible & Mapeo sistemático & Evidencia escasez de estudios experimentales que relacionen arquitectura y consumo energético & Justifica directamente la necesidad de este estudio \\ \hline
\end{tabular}
\end{table}


\section{Descripción de la situación problemática}

\section{Formulación del problema}

\section{Consecuencias de no investigar}

\section{Delimitación}
